% =============================================================================
\section{Señales}
% =============================================================================

\begin{frame}{Señales: definición y clasificación}
  \begin{definicion}
    Una \alert{señal} es una función que transporta información sobre el
    comportamiento de un fenómeno físico: $x\colon \R \to \C$ (tiempo continuo)
    o $x\colon \Zent \to \C$ (tiempo discreto).
  \end{definicion}

  \begin{columns}[T]
    \begin{column}{0.48\textwidth}
      \begin{itemize}
        \item Continua $x(t)$ o discreta $x[n]$
        \item Periódica: $x(t) = x(t + T_0)\ \forall t$
        \item Par / impar: $x(-t) = \pm x(t)$
      \end{itemize}
    \end{column}
    \begin{column}{0.48\textwidth}
      \begin{itemize}
        \item Determinística o aleatoria
        \item Causal: $x(t) = 0$ para $t < 0$
        \item De energía o de potencia
      \end{itemize}
    \end{column}
  \end{columns}
\end{frame}

\begin{frame}{Señales elementales}
  \begin{columns}[c]
    \begin{column}{0.42\textwidth}
      Escalón unitario
      \[
        \escalon{t} =
        \begin{cases}
          1, & t \geq 0\\
          0, & t < 0
        \end{cases}
      \]
      Impulso unitario
      \[
        \impulso{t} = \frac{\dd \escalon{t}}{\dd t},
        \qquad \int_{-\infty}^{\infty} \impulso{t} \dd t = 1
      \]
      \pause
      Exponencial discreta causal
      \[
        x[n] = a^{n}\,\escalon{n}, \quad \abs{a} < 1
      \]
    \end{column}
    \begin{column}{0.54\textwidth}
      \centering
      \begin{tikzpicture}
        \begin{axis}[ejes senal, width=\linewidth, height=3.4cm,
            xlabel={$t$}, ylabel={$\escalon{t}$}, xmin=-2, xmax=3, ymin=0, ymax=1.2,
            xtick={-2,...,3}, ytick={1}]
          \addplot[senal continua, domain=-2:0] {0};
          \addplot[senal continua, domain=0:3] {1};
          \draw[azulutn, very thick, densely dotted] (axis cs:0,0) -- (axis cs:0,1);
        \end{axis}
      \end{tikzpicture}

      \onslide<2->{%
      \begin{tikzpicture}
        \begin{axis}[ejes senal, width=\linewidth, height=3.4cm,
            xlabel={$n$}, ylabel={$x[n]$}, xmin=-3, xmax=10, ymin=0, ymax=1.2,
            xtick={-2,0,...,10}, ytick={1}]
          \addplot[senal discreta, domain=-3:-1, samples=3] {0};
          \addplot[senal discreta, domain=0:10, samples=11] {0.75^x};
        \end{axis}
      \end{tikzpicture}}
    \end{column}
  \end{columns}
\end{frame}

\begin{frame}{Energía y potencia}
  \begin{definicion}[energía y potencia media]
    \vspace{-0.5\baselineskip}
    \begin{align*}
      E_x &= \int_{-\infty}^{\infty} \abs{x(t)}^2 \dd t, &
      P_x &= \lim_{T \to \infty} \frac{1}{2T} \int_{-T}^{T} \abs{x(t)}^2 \dd t
    \end{align*}
  \end{definicion}

  \pause
  \begin{ejemplo}
    Para $x(t) = A\,\e^{-at}\escalon{t}$, con $a > 0$:
    \[
      E_x = \int_{0}^{\infty} A^2 \e^{-2at} \dd t
          = \resaltar<3->{\frac{A^2}{2a}} < \infty
      \quad\Longrightarrow\quad P_x = 0
    \]
    \onslide<3->{Es una \alert{señal de energía}.}
  \end{ejemplo}
\end{frame}
